% Source-derived chapter 04: Springer theory for Hermitian torsion sheaves
% Original source: higher-siegel-weil-unitary-nonsingular-terms
\chapter{Springer theory for Hermitian torsion sheaves}
\label{higher-siegel-weil-unitary-nonsingular-terms-chapter-04}
\source{higher-siegel-weil-unitary-nonsingular-terms}

\begin{remark}[Source section]
\label{higher-siegel-weil-unitary-nonsingular-terms-chapter-04-source}
\source{higher-siegel-weil-unitary-nonsingular-terms}
\source{higher-siegel-weil-unitary-nonsingular-terms}
This chapter reproduces the corresponding section of the retrieved
LaTeX source, including its definitions, statements, examples, and
proofs. It has not yet been translated into Lean.
\notready
\end{remark}

% The source section title was: Springer theory for Hermitian torsion sheaves
\label{sec: Herm}%%%%%
\source{higher-siegel-weil-unitary-nonsingular-terms}
%%%%%%%%%%%%%%%%%%%%%%%%%%%%%%%

In this section we extend the construction in \S\ref{sec: springer} to the case of Hermitian torsion sheaves. The main output is a perverse sheaf $\Spr^{\Herm}_{2d}$ on the moduli stack of Hermitian torsion sheaves with an action of $W_{d}:=(\Z/2\Z)^{d}\rtimes S_{d}$. We will compare the stalks and Frobenius trace functions of $\Spr^{\Herm}_{2d}$ with those of $\Spr_{d}$.


As in \S\ref{sec: springer}, $X$ is a smooth curve over $k$ (not necessarily projective or connected). Recall from \S \ref{ssec: notation} that $\nu: X'\to X$ is a finite map of degree $2$ that is assumed to be {\em generically  \'etale} (and $X'$ is smooth over $k$). We develop the Hermitian  Springer theory in this generality.  Starting from \S\ref{ss:Wd act} we will assume $\nu$ to be \'etale, which is the case needed for proving the main theorem. Let $\s\in\Gal(X'/X)$ be the nontrivial involution. 

\subsection{Local geometry of $\Herm_{d}$}\label{ss:Herm} 
\source{higher-siegel-weil-unitary-nonsingular-terms}
We say that a map of torsion coherent sheaves $a \co \cQ \rightarrow \sigma^* \cQ^\vee$ is \emph{Hermitian} if $\sigma^* a^\vee= a$. 
Let $d\in \N$. Let
\begin{equation*}
\Herm_{d}(X'/X), \textup{ or simply } \Herm_{d}
\end{equation*}
be the moduli stack of pairs $(\cQ,h)$ where $\cQ$ is a torsion coherent sheaf on $X'$ of length $d$, and $h$ is a Hermitian isomorphism $\cQ\isom \s^{*}\cQ^{\vee}:=\s^{*}\un\Ext^{1}(\cQ,\om_{X'})$. 

%We also have a skew-Hermitian variant by changing the condition $\s^{*}h^{\vee}=h$ to $\s^{*}h^{\vee}=-h$. This gives a moduli stack $\SH_{d}(X'/X)$.

We offer two other ways to think about a Hermitian torsion sheaf $(\cQ,h)$. For a torsion sheaf $\cQ$ on $X'$ of length $d$, the datum of a Hermitian isomorphism $h \co \cQ \xrightarrow{\sim} \sigma^* \cQ^\vee$ is equivalent to either:
  \begin{enumerate}
\item a symmetric $k$-bilinear nondegenerate pairing 
$$(\cdot,\cdot): V\times V\to k
$$ on $V=\Gamma(X', \cQ)$ satisfying $(fv_{1},v_{2})=(v_{1},\s^{*}(f)v_{2})$ for any  function $f$ on $X'$ regular near the support of $\cQ$, or 
\item an $\cO_{X'}$-sesquilinear nondegenerate pairing 
$$
\j{\cdot,\cdot}: \cQ\times \cQ\to \om_{F'}/\om_{X'}
$$ satisfying $\j{v_{1},v_{2}}=\s^{*}\j{v_{2},v_{1}}$. Here $\om_{F'}$ is the constant (and quasi-coherent) sheaf on $X'$ whose local sections are the rational $1$-forms on $X'$.
\end{enumerate}
For example, to pass from (2) to (1), form cohomology and apply the trace map $\cohog{0}{\cQ} \rightarrow k$. To pass from $h$ to the pairing in (1), observe that $\ul{\Ext}^1(Q, \omega_{X'}) \cong \ul{\Hom}(Q, \omega_{F'}/\omega_{X'})$ by the long exact sequence associated to $\omega_{X'} \rightarrow  \omega_{F'} \rightarrow \omega_{F'}/\omega_{X'}$. Therefore, $h$ is equivalent to a sesquilinear pairing $\cQ \times \cQ \rightarrow \omega_{F'}/\omega_{X'}$, which upon taking global sections and applying the residue map $\cohog{0}{\omega_{F'}/\omega_{X'}} \rightarrow k$ gives the pairing in (1). 

We refer to $h$, or any of the above equivalent data, as a \emph{Hermitian structure} on $\cQ$. 
  
We have the support map
\begin{equation*}
s^{\Herm}_{d}: \Herm_{d}(X'/X)\to (X'_{d})^{\s}.
% s^{\SH}_{d}: \SH_{d}(X'/X)\to (X'_{d})^{\s}.
\end{equation*}
Note that we have an isomorphism $ (X'_{2d})^{\s} \cong X_d$, sending a $\sigma$-invariant divisor on $X'$ to its descent on $X$. so we will also allow ourselves to view the support map as $s^{\Herm}_{2d}: \Herm_{2d}(X'/X)\to X_d$. 

\begin{remark}
\source{higher-siegel-weil-unitary-nonsingular-terms}
\notready When $\nu$ is \'etale and $d$ is odd, $(X'_{d})^{\s}=\vn$ hence $\Herm_{d}(X'/X)=\vn$. 

In general, when $\nu$ is ramified over the points $R\subset X(\ov k)$, $(X'_{d})^{\s}$ has a decomposition into open and closed subschemes according to the parity of the multiplicities of the divisor at each point $x\in R$.
\end{remark}

Let $\A^{1}_{\sqrt{t}}\to \A^{1}_{t}$ be the square map of affine lines.
\begin{lemma}
\source{higher-siegel-weil-unitary-nonsingular-terms}
\notready\label{l:Herm A1} There is a canonical isomorphism
\source{higher-siegel-weil-unitary-nonsingular-terms}
\begin{eqnarray*}
\Herm_{d}(\A^{1}_{\sqrt{t}}/\A^{1}_{t})\cong [\fro_{d}/\Og_{d}].
%\SH_{d}(\A^{1}_{\sqrt{t}}/\A^{1}_{t})\cong [\sp_{d}/\Sp_{d}].
\end{eqnarray*}
Here $\Og_{d}$ denotes the orthogonal group on a $d$-dimensional nondegenerate quadratic space over $k$ and $\fro_{d}$ is its Lie algebra (the stack $[\fro_{d}/\Og_{d}]$ is independent of the quadratic form).
\end{lemma}
\begin{proof} We give the map $\Herm_{d}(\A^{1}_{\sqrt{t}}/\A^{1}_{t})\to [\fro_{d}/\Og_{d}]$ on $S$-points. For an $S$-point $(\cQ,h)$ of $\Herm_{d}(\A^{1}_{\sqrt{t}}/\A^{1}_{t})$, $V=\Gamma(\A^{1}_{S}, \cQ)$ is a locally free $\cO_{S}$-module of rank $d$ with a nondegenerate symmetric self-duality $(\cdot,\cdot)$, i.e., an $\Og_{d}$-torsor over $S$. Moreover the action of $\sqrt{t}$ on $V$ satisfies $(\sqrt{t}v_{1},v_{2})=-(v_{1},\sqrt{t}v_{2})$ since $\s^{*}\sqrt{t}=-\sqrt{t}$. Therefore $\sqrt{t}$ gives a section of the adjoint bundle of $V$. It is easy to check this map is an equivalence of groupoids $\Herm_{d}(\A^{1}_{\sqrt{t}}/\A^{1}_{t})(S)\isom [\fro_{d}/\Og_{d}](S)$.  
%The proof for $\SH_{d}(\A^{1}_{\sqrt{t}}/\A^{1}_{t})$ is similar.
\end{proof}



An {\em $\s$-equivariant \'etale chart} of $X'_{\ov k}$ is  a pair $(U,f)$, where $U\subset X_{\ov k}$ is an open subset (with preimage $U'\subset X'_{\ov k}$) and a regular function $f: U'\to \A^{1}_{\sqrt{t}, \ov k}$ that is an \'etale map satisfying $\s^{*}f=-f$. Note that if $\nu$ is \'etale, the image of $f$ has to lie in $\A^{1}_{\sqrt{t}, \ov k}\bs\{0\}$.  

A  $\s$-equivariant \'etale chart $(U,f)$ of $X'_{\ov k}$ induces a map 
\[
f^{\Herm}_{d}: \Herm_{d}(U'/U)\to \Herm_{d}(\A^{1}_{\sqrt{t}}/\A^{1}_{t})_{\ov k}
\]
%\[
%f^{\SH}_{d}: \SH_{d}(U'/U)\to \SH_{d}(\A^{1}_{\sqrt{t}}/\A^{1}_{t})_{\ov k}
%\]
by sending $\cQ$ to  $f_{*}\cQ$. Let $\Herm_{d}(U'/U)^{f}$  be the preimages of $(U'_{d})^{\s}\bs \frR^{\s}_{d,f}$ under the support maps (here $\frR_{d,f}\subset U'_{d}$ is defined using the map $f:U'\to \A^{1}_{\sqrt{t},\ov k}$; see \S\ref{ss:Coh loc}). 
 

We have an analog of Lemma \ref{l:Coh local} in the Hermitian setting.

\begin{lemma}
\source{higher-siegel-weil-unitary-nonsingular-terms}
\notready\label{l:Herm local}
\source{higher-siegel-weil-unitary-nonsingular-terms}
\begin{enumerate}
\item Let $(U,f)$ be a $\s$-equivariant \'etale chart for $X'_{\ov k}$. Then the map $f^{\Herm}_{d}$ is \'etale when restricted to $\Herm_{d}(U'/U)^{f}$. 
\item Assume $\nu$ is ramified at at most one point (over $\ov k$). Then the stack $\Herm_{d}(X'/X)_{\ov k}$  is covered by $\Herm_{d}(U'/U)^{f}$ for various $\s$-equivariant \'etale charts $(U,f)$ of $X'_{\ov k}$. In particular,  $\Herm_{d}(X'/X)$ is \'etale locally isomorphic to  $[\fro_{d}/\Og_{d}]$.
\item In general, $\Herm_{d}(X'/X)$ is smooth of dimension $0$.
\end{enumerate}
\end{lemma}
\begin{proof}
(1) is similar to that of Lemma \ref{l:Coh local}(1). 

(2) We only need to construct for $(\cQ,h)\in \Herm_{d}(X'/X)(\ov k)$, a $\s$-equivariant \'etale chart $(U,f)$ such that $(\cQ,h)\in \Herm_{d}(U'/U)^{f}$. Let $Z'$ be its support  in $X'$; since the Hermitian property of $(\cQ, h)$ implies that $Z'$ is stable under $\sigma$, hence is the preimage of some $Z\subset X(\ov k)$ under $\nu$. Let $\cL=(\nu_{*}\cO_{X'_{\ov k}})^{\s=-1}$, a line bundle over $X$. Then the map $r=(r_{z})_{z\in Z}: \cL\to \op_{z\in Z}\cL_{z}/\vp_{z}^{2}$ is surjective. Let $Z_{0}\subset Z$ be the points over which $\nu$ is \'etale (so $Z-Z_{0}$ is empty or has one point). For each $z\in Z_{0}$, upon choosing $z'\in Z'$ over $z$, we may identify $\cL_{z}$ with $\cO_{z}=\ov k[[\vp_{z}]]$; changing $z'$ to $\s(z')$ changes the identification by a sign.  If $z\in Z-Z_{0}$, then $\cL_{z}\cong \sqrt{\vp_{z}}\ov k[[\vp_{z}]]$. Choose a map $c:Z_{0}\to \ov k^{\times}$ such that $c(z)^{2}$ are distinct for $z\in Z$. Let $f$ be a section of $\cL$ over some open neighborhood $U_{1}\subset X_{\ov k}$ of $Z$ such that $r_{z}(f)=c(z)+\vp_{z}$ for $z\in Z_{0}$ under one of the two identifications $\cL_{z}\cong \cO_{z}$, and $r_{z}(f)\equiv\sqrt{\vp_{z}}\mod\vp_{z}$ for $z\in Z-Z_{0}$.  Then $f$ restricts to an \'etale map $U'=\nu^{-1}(U)\to \A^{1}_{\sqrt{t},\ov k}$ for some open neighborhood $U$ of $Z$ in $U_{1}$. The definition of $\cL$ implies $\s^{*}(f)=-f$. Now $\{f(z')|z'\in Z'\}$ is the union of $\{c(z),-c(z)|z\in Z_{0}\}$ and possibly $\{0\}$ if $Z-Z_{0}$ is nonempty, which are all distinct points in $\A^{1}_{\sqrt{t},\ov k}$ by construction. We conclude that $(\cQ,h)\in \Herm_{d}(U'/U)^{f}$. %The argument for $\SH_{d}(X'/X)$ is similar.

(3) Let $R\subset X_{\ov k}$ be the ramification locus of $\nu$. The case $|R|\le 1$ is treated in (2), so we may assume $|R|\ge2$. For $x\in R$, let $Y_{x}=X\bs (R\bs \{x\})$ and let $Y'_{x}=\nu^{-1}(Y_{x})$.   For any function $\d:R\to \Z_{\ge0}$ such that $\sum_{x\in R} \d(x)=d$ we have a map $\frY_{\d}:=\prod_{x\in R}(Y'_{x,d(x)})^{\s}\to (X'_{d})^{\s}$ by adding divisors. Let $\frY_{\d}^{\hs}\subset \frY_{\d}$ be the open locus where the divisors indexed by different $x\in R$ are disjoint. It is clear that $\frY^{\hs}_{\d}\to (X'_{d})^{\s}$ is \'etale and for varying $\d$ their images cover $(X'_{d})^{\s}$. To prove the statement it suffices to show that the base change $\Herm_{d}(X'/X)|_{\frY_{\d}^{\hs}}$ is smooth of dimension $0$ for each $\d$. Observe that $\Herm_{d}(X'/X)|_{\frY_{\d}^{\hs}}$ is isomorphic to the restriction of the product $\prod_{x\in R}\Herm_{\d(x)}(Y'_{x}/Y_{x})$ to $\frY_{\d}^{\hs}$. Since $\nu|_{Y_{x}}: Y'_{x}\to Y_{x}$ is ramified at one point, by (2) $\Herm_{\d(x)}(Y'_{x}/Y_{x})$ is smooth of dimension $0$. Therefore $\Herm_{d}(X'/X)|_{\frY_{\d}^{\hs}}\cong \prod_{x\in R}\Herm_{\d(x)}(Y'_{x}/Y_{x})|_{\frY_{\d}^{\hs}}$ is smooth of dimension $0$.  
\end{proof}

%\begin{remark} When $\nu$ is ramified, the $\Herm_{d}(U'/U)^{f}$ for $\s$-equivariant \'etale charts $(U,f)$  only cover the locus of those $(\cQ,h)$ whose support contains at most one ramified point of $\nu$.
%\end{remark}

\begin{remark}
\source{higher-siegel-weil-unitary-nonsingular-terms}
\notready
There is an obvious notion of skew-Hermitian torsion sheaves. Let $\SH_{d}(X'/X)$ be the moduli stack of skew-Hermitian torsion sheaves on $(X',\s)$ of length $d$. Then $d$ is even if $\SH_{d}(X'/X)\ne\vn$. The skew-Hermitian analog of Lemma \ref{l:Herm local} says that $\SH_{d}(X'/X)$ is \'etale locally isomorphic to $[\sp_{d}/\Sp_{d}]$, at least  when $\nu$ is \'etale.
\end{remark}



\subsection{The Hermitian Springer sheaf}\label{ss:Herm Spr}  
\source{higher-siegel-weil-unitary-nonsingular-terms}
Let $\wt\Herm_{d}(X'/X)$ be the moduli stack classifying $(\cQ,h)\in \Herm_{d}(X'/X)$ together with a full flag
\begin{equation*}
0\subset \cQ_{1}\subset \cdots \subset \cQ_{i}\subset \cdots \subset \cQ_{d-1}=\cQ_{1}^{\bot}\subset \cQ_{d}=\cQ,
\end{equation*} 
where $\cQ_{i}$ has length $i$ and $\cQ_{d-i}=\cQ_{i}^{\bot}$ (the orthogonal of $\cQ_{i}$ under the Hermitian pairing $\cQ\times \cQ\to \om_{F'}/\om_{X'}$). Let
\begin{equation*}
\pi^{\Herm}_{d}: \wt\Herm_{d}(X'/X)\to \Herm_{d}(X'/X)
\end{equation*}
be the forgetful map.  Let $\Herm_{d}(X'/X)^{\c}\subset \Herm_{d}(X'/X)$ be the preimage of the multiplicity-free part  $X'^{\c}_{d}$ under the support map $s^{\Herm}_{d}$.

%Similarly we define $\pi^{\SH}_{d}: \wt\SH_{d}(X'/X)\to \SH_{d}(X'/X)$.


We recall the Grothendieck alteration for the full orthogonal group $\Og(V,Q)$ for some vector space $V$ of dimension $d$ over $k$ and a nondegenerate quadratic form $Q$ on $V$. Let $\Fl(V,Q)$ be the flag variety that parametrizes full isotropic flags $V_{\bu}=(V_{1}\subset \cdots\subset V_{d}=V)$ in $V$. Note that when $d$ is even, this is different from the flag variety of $\SO(V,Q)$ but rather a double cover of it because there are two choices for $V_{d/2}$ given the rest of members of a flag. Let $\fro(V,Q)$ be the Lie algebra of $\Og(V,Q)$. Let $\wt\fro(V,Q)$ be the moduli space of pairs $(A, V_{\bu})\in \fro(V,Q)\times\Fl(V,Q)$ such that $AV_{i}\subset V_{i}$ for all $i$. The Grothendieck alteration for $\Og(V,Q)$ is the $\Og(V,Q)$-equivariant map $\wt\fro(V,Q)\to \fro(V,Q)$ forgetting the flag. The quotient stacks $[\fro(V,Q)/\Og(V,Q)]$ and $[\wt\fro(V,Q)/\Og(V,Q)]$ are canonical; they are independent of the quadratic form $Q$ and only depend on $d=\dim V$. Therefore we also write the Grothendieck alteration as $\pi_{\Og_{d}}: [\wt\fro_{d}/\Og_{d}]\to [\fro_{d}/\Og_{d}]$.


\begin{prop}
\source{higher-siegel-weil-unitary-nonsingular-terms}
\notready\label{p:SH Spr}
\source{higher-siegel-weil-unitary-nonsingular-terms}
\begin{enumerate}
\item If $\nu$ is ramified at at most one point, then the map $\pi^{\Herm}_{d}$ is \'etale locally isomorphic to the Grothendieck alteration $\pi_{\Og_{d}}:[\wt\fro_{d}/\Og_{d}]\to [\fro_{d}/\Og_{d}]$. 
\item In general, $\wt\Herm_{d}(X'/X)$ is smooth of dimension $0$ and $\pi^{\Herm}_{d}$ is a small map. In particular, the complex
\begin{equation*}
\Spr^{\Herm}_{d}:=R\pi^{\Herm}_{d,*}\Qlbar
\end{equation*}
is the middle extension perverse sheaf of its restriction to $\Herm_{d}(X'/X)^{\c}$.
\end{enumerate}
%\item The map $\pi^{\SH}_{2d}$ is \'etale locally isomorphic to the Grothendieck alteration $[\wt\sp_{2d}/\Sp_{2d}]\to [\sp_{2d}/\Sp_{2d}]$. In particular, it is small and
%\begin{equation}
%\Spr^{\SH}_{2d}:=R\pi^{\SH}_{2d,*}\Qlbar
%\end{equation}
%is the middle extension perverse sheaf of its restriction to $\SH_{2d}(X'/X)^{\c}$, and carries a canonical $W_{d}$-action. 
\end{prop}
\begin{proof} 
(1) The proof is similar to that of Corollary \ref{c:Coh Spr}. For a $\s$-equivariant \'etale chart $(U,f)$ for $X'_{\ov k}$ we have a diagram with Cartesian squares and \'etale horizontal maps by Lemma \ref{l:Herm local}
\begin{equation}\label{SH chart}
\source{higher-siegel-weil-unitary-nonsingular-terms}
\xymatrix{    \wt\Herm_{d}(X'/X)_{\ov k} \ar[d]^{\pi^{\Herm}_{d,X'/X}}& \wt\Herm_{d}(U'/U)^{f} \ar[r]\ar@{_{(}->}[l] \ar[d]^{\pi^{\Herm}_{d,U'/U}}& \wt\Herm_{d}(\A^{1}_{\sqrt{t}}/\A^{1}_{t}\ar[d]^{\pi^{\Herm}_{d, \A^{1}_{\sqrt{t}}/\A^{1}_{t}}})_{\ov k}        \\
\Herm_{d}(X'/X)_{\ov k}    & \Herm_{d}(U'/U)^{f}\ar[r]^{f^{\Herm}_{d}}\ar@{_{(}->}[l] & \Herm_{d}(\A^{1}_{\sqrt{t}}/\A^{1}_{t})_{\ov k}     }
\end{equation}
Using the isomorphism in Lemma \ref{l:Herm A1}, we identify $\pi^{\Herm}_{d, \A^{1}_{\sqrt{t}}/\A^{1}_{t}}$ with the Grothendieck alteration $\pi_{\Og_{d}}$.  Since $\Herm_{d}(U'/U)^{f}$ cover $\Herm_{d}(X'/X)$ by Lemma \ref{l:Herm local}(2), $\pi^{\Herm}_{d,X'/X}$ is 
\'etale locally  isomorphic to $\pi^{\Herm}_{d, \A^{1}_{\sqrt{t}}/\A^{1}_{t}}=\pi_{\Og_{d}}$.

(2) We use the notation from the proof of Lemma \ref{l:Herm local}(3). We may assume $|R|\ge2$. For each function $\d: R\to \Z_{\ge0}$ satisfying $\sum_{x\in R}\d(x)=d$, the base change of $\pi^{\Herm}_{d}$ along $\frY_{\d}^{\hs}\to (X'_{d})^{\s}$ is a disjoint union of the restriction of
\begin{equation*}
\prod_{x}\pi^{\Herm}_{\d(x)}: \prod_{x\in R}\wt\Herm_{\d(x)}(Y'_{x}/Y_{x})\to \prod_{x\in R}\Herm_{\d(x)}(Y'_{x}/Y_{x})
\end{equation*}
to $\frY_{\d}^{\hs}$. The disjoint union comes from different ways to distribute $\supp(\cQ_{i}/\cQ_{i-1})$ among various factors in the product $\prod_{x}(Y'_{x,\d(x)})^{\s}$. By (1), $\pi^{\Herm}_{\d(x)}: \wt\Herm_{\d(x)}(Y'_{x}/Y_{x})\to \Herm_{\d(x)}(Y'_{x}/Y_{x})$ has smooth $0$-dimensional source and is small for each $x\in R$, the same holds true for the base change of $\pi^{\Herm}_{d}$ to $\frY_{\d}^{\hs}$. Since $\{\frY_{\d}^{\hs}\}_{\d}$ form an \'etale covering of $(X'_{d})^{\s}$, the same is true for $\pi^{\Herm}_{d}$.
\end{proof}


\subsection{The action of $W_{d}$} \label{ss:Wd act} 
\source{higher-siegel-weil-unitary-nonsingular-terms}
From now on we assume that $\nu:X'\to X$ is an \'etale double cover.  In this case, $(X'_{2d})^{\s}$ can be identified with $X_{d}$ via $\nu^{-1}(D)\bij D$. Let
\begin{equation*}
W_{d}:=(\Z/2\Z)^{d}\rtimes S_{d}
\end{equation*}
be the Weyl group for $\Og_{2d}$. Then $\pi^{\Herm}_{2d}$ is a $W_{d}$-torsor over $\Herm_{2d}(X'/X)^{\c}$.  

\begin{cor}[of Proposition \ref{p:SH Spr}(1)]
\source{higher-siegel-weil-unitary-nonsingular-terms}
\notready\label{c:Wd action} If $\nu:X'\to X$ is an \'etale double cover, then there is a canonical action of $W_{d}$ on $\Spr^{\Herm}_{2d}$ extending the geometric action on its restriction to $\Herm_{2d}(X'/X)^{\c}$.
\source{higher-siegel-weil-unitary-nonsingular-terms}
\end{cor}


\begin{defn}
\source{higher-siegel-weil-unitary-nonsingular-terms}
\notready\label{d:Spr isotypic} 
\source{higher-siegel-weil-unitary-nonsingular-terms}
\begin{enumerate}
\item For any representation $\r$ of $W_{d}$, we define $\Spr^{\Herm}_{2d}[\r]$ to be the perverse sheaf  on $\Herm_{2d}(X'/X)$:
\begin{equation*}
\Spr^{\Herm}_{2d}[\r]=(\r^{\vee}\ot \Spr^{\Herm}_{2d})^{W_{d}}\in D^b(\Herm_{2d}(X'/X),\Qlbar).
\end{equation*}
\item We define the Hermitian  analog of the Springer sheaf $\Spr$ as
\begin{eqnarray*}
\HSpr_{d}:=(\Spr^{\Herm}_{2d})^{(\Z/2\Z)^{d}}\in D^b(\Herm_{2d}(X'/X),\Qlbar).
%\SHSpr_{d}:=(\Spr^{\SH}_{2d})^{(\Z/2\Z)^{d}}\in \Perv(\SH_{2d}(X'/X)).
\end{eqnarray*}
\end{enumerate}
\end{defn}
Note that the notation shifts from the subscript $2d$ to $d$. By Corollary \ref{c:Wd action}, $\HSpr_{d}$ carries a canonical $S_{d}$-action.

\begin{remark}
\source{higher-siegel-weil-unitary-nonsingular-terms}
\notready In the case $\nu$ is ramified with  ramification locus $R\subset X(\ov k)$, the stack $\Herm_{2d}(X'/X)_{\ov k}$ decomposes into the disjoint union of open and closed substacks $\Herm^{\ep}_{2d}(X'/X)_{\ov k}$ indexed by $\e: R\to \{0,1\}$ where the length of $\cQ_{x}$ has parity $\e(x)$ for all $x\in R$. Then $\Spr^{\Herm}_{2d}|_{\Herm^{\ep}_{2d}(X'/X)_{\ov k}}$ carries a canonical action of $W_{d'}$ where $d'=(d-\sum_{x\in R}\ep(x))/2$.  
\end{remark}

%%%% above ok %%%%%

\subsection{The Springer fibers over $\Herm_{2d}$}\label{ss:Herm fiber} 
\source{higher-siegel-weil-unitary-nonsingular-terms}
Let $(\cQ,h)\in \Herm_{2d}(X'/X)(\ov k)$ and consider its Hermitian Springer fiber
\begin{equation*}
\cB^{\Herm}_{\cQ}:=\pi^{\Herm,-1}_{2d}(\cQ,h).
\end{equation*}
This is a proper scheme over $\ov k$. In this subsection we prove the Hermitian analogs of results in \S\ref{ss:Spr fiber}.

Let $Z'=\supp\cQ\subset X'(\ov k)$.  Let $D'=s^{\Herm}_{2d}(\cQ)\in (X'_{2d})^{\s}(\ov k)$, which is of the form $D'=\nu^{-1}(D)$ for some $D\in X_{d}(\ov k)$. Let $Z=\nu(Z')$, the support of $D$. Write $D'=\sum_{z\in Z'}d_{z}z$. 

Let $\Sigma(Z')$ be the set of maps $y': \{1,2,\cdots,2d\}\to Z'$ satisfying $y'(2d+1-i)=\s(y'(i))$ for all $i$ and $\sum_{i=1}^{2d}y'(i)=D'$. Identifying $W_{d}$ with permutations of $\{1,2,\cdots,2d\}$ commuting with the involution $i\mapsto 2d+1-i$, we get an action of $W_{d}$ on $\Sigma(Z')$ by $w: y'\mapsto y'\c w^{-1}$.

Similarly let $\Sigma(Z)$ be the set of maps $y:\{1,\cdots, d\}\to Z$ such that $\sum_{i=1}^{d}y(i)=D$. Then the natural map $\Sigma(Z')\to \Sigma(Z)$ (sending $y'$ to $y$ defined by $y(i)=\nu(y'(i))$) is a $(\Z/2\Z)^{d}$-torsor. 

For $y'\in\Sigma(Z')$, let $\cB^{\Herm}_{\cQ}(y')$ be the subscheme of $\cB^{\Herm}_{\cQ}$ consisting of isotropic flags $\cQ_{\bu}$ such that $\supp(\cQ_{i}/\cQ_{i-1})=y'(i)$ for all $1\le i\le 2d$. Then we have a decomposition into open and closed subschemes
\begin{equation*}
\cB^{\Herm}_{\cQ}=\coprod_{y'\in \Sigma(Z')}\cB^{\Herm}_{\cQ}(y').
\end{equation*}
Accordingly we get a decomposition of cohomology
\begin{equation*}
\cohog{*}{\cB^{\Herm}_{\cQ}}=\bigoplus_{y'\in \Sigma(Z')}\cohog{*}{\cB^{\Herm}_{\cQ}(y')}.
\end{equation*}


\begin{lemma}
\source{higher-siegel-weil-unitary-nonsingular-terms}
\notready\label{l:permute comp}
\source{higher-siegel-weil-unitary-nonsingular-terms}
The action of $w\in W_{d}$ on $\cohog{*}{\cB^{\Herm}_{\cQ}}$ sends the direct summand $\cohog{*}{\cB^{\Herm}_{\cQ}(y')}$ to the direct summand $\cohog{*}{\cB^{\Herm}_{\cQ}(y'\c w^{-1})}$.
\end{lemma}
\begin{proof}
It suffices to check the statement for each simple reflection $s_{i}$, $i=1,\cdots, d$. Here, for $1\le i\le d-1$,  $s_{i}=(i,i+1)(2d-i,2d+1-i)$; for $i=d$, $s_{d}=(d,d+1)$. For $1\le i\le d$, let $\wt\Herm^{i}_{2d}$ be the moduli stack classifying isotropic flags that only misses the terms $\cQ_{i}$ and $\cQ_{2d-i}$ (for $i=d$ only misses $\cQ_{d}$). Then we have a factorization
\begin{equation*}
\pi^{\Herm}_{2d}: \wt\Herm_{2d}\xr{\r_{i}}\wt\Herm^{i}_{2d}\xr{\pi_{i}}\Herm_{2d}.
\end{equation*}
The map $\r_{i}$ is an \'etale double cover over the open dense locus  $\wt\Herm^{i,\hs}_{2d}$  where $\cQ_{i+1}/\cQ_{i-1}$ (which has length $2$) is supported at two distinct points. The map $\r_{i}$ is small, and $R\r_{i*}\Qlbar$ carries an involution $\wt s_{i}$, which induces an involution $\wt s_{i}$ on $R\pi_{i*}R\r_{i*}\Qlbar\cong R\pi^{\Herm}_{2d*}\Qlbar$. This action coincides with the action of $s_{i}$ over $\Herm_{2d}^{\c}$, hence coincides with $s_{i}$ everywhere. 

Let $(\cQ,h)\in\Herm_{2d}(\ov k)$, and $\cB_{\cQ}^{i}=\pi_{i}^{-1}(\cQ,h)$. We have a decomposition of $\cB_{\cQ}^{i}$ by the orbit set $\Sigma(Z')/\j{s_{i}}$. When $y'\in\Sigma(Z')$ satisfies $y'\ne y'\c s_{i}$, the $s_{i}$-orbit of $\y'=\{y',y'\c s_{i}\}$ gives an open closed substack $\cB^{i}_{\cQ}(\y')\subset \cB_{\cQ}^{i}$, such that $\r^{-1}_{i}(\cB^{i}_{\cQ}(\y'))=\cB_{\cQ}(y')\coprod\cB_{\cQ}(y'\c s_{i})$, and $\cB^{i}_{\cQ}(\y')\subset \wt\Herm^{i,\hs}_{2d}$. Therefore in this case the action of $\wt s_{i}$ on $\cohog{*}{\r^{-1}_{i}(\cB^{i}_{\cQ}(\y'))}$ comes from the involution on $\cB_{\cQ}(y')\coprod\cB_{\cQ}(y'\c s_{i})$ that interchanges the two components. Since $\wt s_{i}=s_{i}$, this proves the statement for $s_{i}$ and $y'$ such that $y'\ne y'\c s_{i}$. For $y'=y'\c s_{i}$ the statement is vacuous. This finishes the proof. 
\end{proof}






Choose $Z^{\sh}\subset Z'$ such that $Z^{\sh}\coprod\s(Z^{\sh})=Z'$. Then for each $x\in Z$ there is a unique $x^{\sh}\in Z^{\sh}$ above $x$. For  $y'\in \Sig(Z')$ with image $y\in \Sig(Z)$, we have an isomorphism
\begin{equation}\label{gZ}
\source{higher-siegel-weil-unitary-nonsingular-terms}
\g_{Z^{\sh}, y'}: \cB^{\Herm}_{\cQ}(y')\isom\cB_{\nu_{*}(\cQ|_{Z^{\sh}})}(y)\cong \prod_{x\in Z}\cB_{\cQ_{x^{\sh}}}
\end{equation}
mapping  $(\cQ_{i})_{1\le i\le 2d}$ to the (non-strictly increasing) flag  $(\cQ_{i,x^{\sh}})$ of $\cQ_{x^{\sh}}$.

If $y',y''\in \Sig(Z')$, the composition
\begin{equation*}
\g_{y',y''}:=\g_{Z^{\sh},y''}^{-1}\c\g_{Z^{\sh}, y'}:\cB^{\Herm}_{\cQ}(y')\isom\cB^{\Herm}_{\cQ}(y'')
\end{equation*}
is independent of the choice of $Z^{\sh}$.

\begin{lemma}
\source{higher-siegel-weil-unitary-nonsingular-terms}
\notready\label{l:min w Herm}
\source{higher-siegel-weil-unitary-nonsingular-terms}
Let $y',y''\in \Sig(Z')$ and let $w\in W_{d}$ be a minimal length element such that $y''=y'\c w^{-1}$. Then the Springer action $w: \cohog{*}{\cB^{\Herm}_{\cQ}(y')}\to  \cohog{*}{\cB^{\Herm}_{\cQ}(y'')}$ is induced by the isomorphism $\g_{y', y''}$.
\end{lemma}
\begin{proof}
Similar to the proof of Lemma \ref{l:min w}.
\end{proof}



\subsection{Comparing stalks of $\HSpr_{d}$ and $\Spr_{d}$}\label{ss:HSpr stalk}
\source{higher-siegel-weil-unitary-nonsingular-terms}
In this subsection we abbreviate $\Herm_{2d}(X'/X)$ by $\Herm_{2d}$. Consider the stack $\Lagr_{2d}$ classifying pairs $(\cL\subset\cQ)$ where $\cQ\in\Herm_{2d}$ and $\cL\subset \cQ$ is a Lagrangian subsheaf, i.e., $\cL$ has length $d$ and the composition $\cL\inj \cQ\xr{h}\s^{*}\cQ^{\vee}\to \s^{*}\cL^{\vee}$ is zero. We have natural maps
\begin{equation*}
\xymatrix{\Herm_{2d} & \Lagr_{2d} \ar[l]_{\upsilon_{2d}}\ar[r]^{\e'_{d}}& \Coh_{d}(X')\ar[r]^{\nu_{*}} & \Coh_{d}(X)}
\end{equation*}
where $\upsilon_{2d}(\cL\subset\cQ)=\cQ$ and $\e'_{d}(\cL\subset\cQ)=\cL$.  Let $\e_{d}=\nu_{*}\c \e'_{d}: \Lagr_{2d}\to \Coh_{d}(X)$.
 
Let $(X'_{d})^{\dm}\subset X'_{d}$ be the open subscheme parametrizing $D\in X'_{d}$ such that $D\cap\s(D)=\vn$.  
Let $\Lagr^{\dm}_{2d}\subset \Lagr_{2d}$ be the preimage of $(X'_{d})^{\dm}$ under the map $\Lagr_{2d}\xr{\e'_{d}}\Coh_{d}(X')\xr{s^{\Coh}_{d,X'}} X'_{d}$.  It is easy to see that $\e'_{d}$ restricts to an isomorphism
\begin{equation*}
\Lagr^{\dm}_{2d}\cong \Coh_{d}(X')^{\dm}
\end{equation*}
whose inverse is given by $\cL\mapsto (\cL\subset \cQ=\cL\op \s^{*}\cL^{\vee})$. 

Let $\upsilon^{\dm}_{2d}$ and $\e^{\dm}_{d}$ be the restrictions of $\upsilon_{2d}$ and $\e_{d}$ to $\Lagr^{\dm}_{2d}$. Thus we view $\Lagr^{\dm}_{2d}\cong \Coh_{d}(X')^{\dm}$ as a correspondence between $\Herm_{2d}$ and $\Coh_{d}(X)$
\begin{equation*}
\xymatrix{ \Herm_{2d} &  \Lagr^{\dm}_{2d}\ar[l]_{\upsilon^{\dm}_{2d}}\ar[r]^{\e^{\dm}_{d}} &\Coh_{d}(X) }
\end{equation*}
Note that both $\upsilon^{\dm}_{2d}$ and $\e^{\dm}_{d}$ are surjective. Now both $\Herm_{2d}$ and $\Coh_{d}(X)$ carry Springer sheaves $\HSpr_{d}$ and $\Spr_{d}$ with $S_{d}$-actions. The following proposition says that they become isomorphic after pullback to $\Lagr^{\dm}_{2d}$.




\begin{prop}
\source{higher-siegel-weil-unitary-nonsingular-terms}
\notready\label{p:Spr HSpr}
\source{higher-siegel-weil-unitary-nonsingular-terms}
There is a canonical $S_{d}$-equivariant isomorphism of perverse sheaves on $\Lagr^{\dm}_{2d}$
\begin{equation*}
\upsilon^{\dm,*}_{2d}\HSpr_{d}\cong \e^{\dm,*}_{d}\Spr_{d}.
\end{equation*}
\end{prop}
\begin{proof} The map $\pi^{\Herm}_{2d}$ factors as
\begin{equation*}
\pi^{\Herm}_{2d}: \wt\Herm_{2d}\xr{\l_{2d}} \Lagr_{2d}\xr{\upsilon_{2d}} \Herm_{2d}.
\end{equation*}
Let $\l^{\dm}_{2d}: \wt\Herm^{\dm}_{2d}\to \Lagr^{\dm}_{2d}$ be the restriction of $\l_{2d}$ to $\Lagr^{\dm}_{2d}$.
We have a commutative diagram
\begin{equation}\label{3sq}
\source{higher-siegel-weil-unitary-nonsingular-terms}
\xymatrix{\Herm_{2d}\times_{X_{d}}X^{d}\ar[d] & \wt\Herm^{\dm}_{2d}\ar[l]\ar[dl]_{\pi^{\Herm,\dm}_{2d}}\ar[d]^{\l^{\dm}_{2d}}\ar[r]^-{\wt \e'^{\dm}_{d}} & \wt\Coh_{d}(X')^{\dm}\ar[d]^{\pi^{\Coh,\dm}_{X',d}}\ar[r]^{\nu_{*}} &\wt\Coh_{d}(X)\ar[d]^{\pi^{\Coh}_{d}} \\
\Herm_{2d} & \Lagr^{\dm}_{2d}\ar@/_2pc/[rr]^-{\e^{\dm}_{d}}\ar[l]_{\upsilon^{\dm}_{2d}}\ar[r]^-{\e'^{\dm}_{d}} & \Coh_{d}(X')^{\dm}\ar[r]^{\nu_{*}} &   \Coh_{d}(X)
}
\end{equation}
Here $\Coh_{d}(X')^{\dm}$ and $\wt\Coh_{d}(X')^{\dm}$ are the preimages of $(X'_{d})^{\dm}$ under the support map.  We have:
\begin{itemize}
\item The middle square is Cartesian. This is true even before restricting to the $\dm$ locus. 
\item Since $\e'^{\dm}_{d}$  is an isomorphism, so is $\wt \e'^{\dm}_{d}$.
\item The rightmost square is Cartesian. 
\end{itemize}
From these properties we get maps
\begin{eqnarray*}
\a: \upsilon^{\dm *}_{2d}\HSpr_{d} &\to& \upsilon^{\dm*}_{2d}\Spr^{\Herm}_{2d}=\upsilon^{\dm*}_{2d}R\upsilon_{2d*}R\l_{2d*}\Qlbar\to \upsilon^{\dm*}_{2d}R\upsilon^{\dm}_{2d*}R\l^{\dm}_{2d*}\Qlbar \\
&\to&  R\l^{\dm}_{2d*}\Qlbar\cong \e^{\dm*}_{d}\pi^{\Coh}_{d*}\Qlbar=\e^{\dm*}_{d}\Spr_{d}.
\end{eqnarray*}



To check $\a$ is an isomorphism, it suffices to check on geometric stalks. Let $\cL\in \Coh_{d}(X')^{\dm}(\ov k)$ with support $Z^{\sh}\subset X'(\ov k)$. Let $\cQ=\cL\op \s^{*}\cL^{\vee}\in \Herm_{2d}(\ov k)$. The support of $\cQ$ is $Z'=Z^{\sh}\coprod \s(Z^{\sh})$, with image $Z\subset X(\ov k)$. We have $(\cL\subset \cQ)\in \Lagr^{\dm}_{2d}(\ov k)$, with image $\nu_{*}\cL\in\Coh_{d}(X)(\ov k)$. The stalk of $\a$ at $(\cL\subset\cQ)$ is
\begin{equation}\label{astalk}
\source{higher-siegel-weil-unitary-nonsingular-terms}
\a_{(\cL\subset\cQ)}: \cohog{*}{\cB^{\Herm}_{\cQ}}^{(\Z/2\Z)^{d}}\to\cohog{*}{\l^{-1}_{2d}(\cL\subset\cQ)}\isom \cohog{*}{\cB_{\nu_{*}\cL}}
\end{equation}
Recall $\cB^{\Herm}_{\cQ}=\coprod_{y'\in\Sig(Z')}\cB^{\Herm}_{\cQ}(y')$. Let $\Sig(Z^{\sh})\subset \Sig(Z')$ be the set of $y'$ such that $y'(i)\in Z^{\sh}$ for $1\le i\le d$. Then we have a natural bijection $\Sig(Z)\bij\Sig(Z^{\sh}), y\mapsto y^{\sh}$. The fiber $\l^{-1}_{2d}(\cL\subset\cQ)$ is the disjoint union $\coprod_{y\in\Sig(Z)}\cB^{\Herm}_{\cQ}(y^{\sh})$. Recall the isomorphism $\g_{Z^{\sh}, y^{\sh}}:  \cB^{\Herm}_{\cQ}(y^{\sh})\isom  \cB_{\nu_{*}\cL}(y)$ from \eqref{gZ}. Using these descriptions we may rewrite \eqref{astalk} as
\begin{equation}
\cohog{*}{\cB^{\Herm}_{\cQ}}^{(\Z/2\Z)^{d}}  \surj \op_{y'\in\Sig(Z^{\sh})}\cohog{*}{\cB^{\Herm}_{\cQ}(y')}\cong \op_{y\in \Sig(Z)}\cohog{*}{\cB_{\nu_{*}\cL}(y)}=\cohog{*}{\cB_{\nu_{*}\cL}}.
\end{equation}
It remains to show that the first map above is an isomorphism. But this follows from the fact that  $(\Z/2\Z)^{d}$ acts freely on $\Sig(Z')$ with orbit representatives $\Sig(Z^{\sh})$, and Lemma \ref{l:permute comp}. This shows that $\a$ is an isomorphism. 


Finally we show that $\a$ is $S_{d}$-equivariant. By Proposition  \ref{p:SH Spr}, $\pi^{\Herm}_{2d}$ and hence $\l_{2d}$ is small,  $R\l_{2d*}\Qlbar$ is the middle extension from a dense open substack of $\Lagr_{2d}$. Therefore the same is true for $R\l^{\dm}_{2d*}\Qlbar$. Since $\a$ is an isomorphism, both $\upsilon^{\dm,*}_{2d}\HSpr_{d}$ and $\e^{\dm*}_{d}\Spr_{d}$ are middle extension perverse sheaves from a dense open substack of $\Lagr^{\dm}_{2d}$. To check that $\a$ is $S_{d}$-equivariant it suffices to check it over the dense open substack which is the preimage of the multiplicity-free locus $X_{d}^{\c}$. Over $X_{d}^{\c}$, all squares in \eqref{3sq} are Cartesian,  and all vertical maps are $S_{d}$-torsors. The $S_{d}$-actions on $\HSpr_{d}|_{\Herm_{2d}^{\c}}$ and $\Spr_{d}|_{\Coh_{d}(X)^{\c}}$ come from the vertical $S_{d}$-torsors in the diagram, so $\a$ is $S_{d}$-equivariant when restricted over $X_{d}^{\c}$. This finishes the proof.
\end{proof}

%There is an obvious analogue of this proposition for $\SH_{2d}$ instead of $\Herm_{2d}$.

%\begin{cor}\label{c:Wind}
%Let $y\in \Sig(Z)$. There is an $S_{d}$-equivariant isomorphism
%\begin{equation}
%(\HSpr_{d})_{\cQ}\cong \Ind^{S_{d}}_{S_{y}}\left(\bigotimes_{x\in Z}(\HSpr_{d_{x}})_{\cQ_{\nu^{-1}(x)}}\right).
%\end{equation}
%Here $S_{y}\cong \prod_{x\in Z}S_{I_{x}}$ is the stabilizer of $y$ under $S_{d}$ (where $I_{x}\subset\{1,2,\cdots, d\}$ is the preimage of $x$ under $y$, $d_{x}=\#I_{x}$), and the factor $S_{I_{x}}$ acts on the tensor factor $(\HSpr_{d_{x}})_{\cQ_{\nu^{-1}(x)}}$.
%\end{cor}
%\begin{proof}
%Let $\Sig(Z')_{y}\subset \Sig(Z)$ be the preimage of $y$. Let $W_{y}$ be the preimage of $S_{y}$ in $W_{d}$. By Lemma \ref{l:permute comp}, $\op_{y'\in \Sig(Z')_{y}} \cohog{*}{\cB^{\Herm}_{\cQ}(y')}$ is stable under $W_{y}$, and it induced a canonical isomorphism
%\begin{equation}
%\cohog{*}{\cB^{\Herm}_{\cQ}}\cong \Ind_{W_{y}}^{W_{d}}(\op_{y'\in \Sig(Z')_{y}} \cohog{*}{\cB^{\Herm}_{\cQ}(y')}).
%\end{equation}
%On the other hand, $\coprod_{y'\in \Sig(Z')_{y}}\cB^{\Herm}_{\cQ}(y')\cong \prod_{x\in Z}\cB^{\Herm}_{\cQ_{\nu^{-1}(x)}}$. Write $W_{y}=\prod_{x\in Z}W_{I_{x}}$, then $W_{I_{x}}$ acts on each factor 
%\end{proof}
%

%In particular, $\pi^{\Lagr}_{2d}$ is small, and $R\pi^{\Lagr}_{2d*}\Qlbar$ carries an action of $S_{d}$. After taking $R\upsilon^{2d*}$ this action is the restriction of the $W_{d}$ action on $\Spr^{\Herm}_{2d}$. On the other hand, the $S_{d}$-action on $R\pi^{\Lagr}_{2d*}\Qlbar$ is induced from the $S_{d}$-action on $\Spr_{X',d}=R\pi^{\Coh}_{X', d*}\Qlbar$ constructed in  Corollary \ref{c:Coh Spr}.

%Let $e^{\dm}_{d}$ be the restriction of $e_{d}$  to $\Lagr^{\dm}_{2d}$, then $e^{\dm}_{d}: \Lagr^{\dm}_{2d}\incl \Coh_{d}(X')$ is an open embedding with image $\Coh_{d}(X')^{\dm}$ equal to the preimage of $(X'_{d})^{\dm}$ under $s^{\Coh}_{X',d}: \Coh_{d}(X')\to (X'_{d})^{\dm}$).  The inverse $\Coh_{d}(X')^{\dm}\isom \Lagr^{\dm}_{2d}$ is given by $\cL\mapsto \cL\op\s^{*}\cL^{\vee}$ with the natural Hermitian pairing. Let $\e^{\dm}_{d}$ be the composition 
%\begin{equation}
%\e^{\dm}_{d}: \Lagr^{\dm}_{2d}\xr{e^{\dm}_{d}}\Coh_{d}(X')^{\dm}\xr{\nu_{*}}\Coh_{d}(X).
%\end{equation}

%We have a Cartesian diagram
%\begin{equation}
%\xymatrix{\Coh_{d}(X'/X)^{\dm}\ar[r]^-{s^{\Coh}_{X',d}}\ar[d]^{\nu_{*}} & (X'_{d})^{\dm}\ar[d]^{\nu_{d}}\\
%\Coh_{d}(X)\ar[r]^{s_{d}^{\Coh}} & X_{d}}
%\end{equation}



%Let $\Sig(Z^{\sh})\subset \Sig(Z')$ be the set of $y'$ such that $y'(i)\in Z^{\sh}$ for $1\le i\le d$. Then we have a natural bijection $\Sig(Z^{\sh})\isom\Sig(Z)$, and $\Sig(Z^{\sh})$ contains exactly one element from each $(\Z/2\Z)^{d}$-orbit on $\Sig(Z')$. We have
%\begin{equation}
%\pi^{\Lagr,-1}_{2d}(\cQ^{\sh})=\coprod_{y'\in \Sig(Z^{\sh})}\cB^{\Herm}_{\cQ}(y')\subset \cB^{\Herm}_{\cQ}.
%\end{equation}
%By the Cartesian diagram \eqref{LagCoh}, we have a canonical isomorphism
%\begin{equation}\label{QshQfl}
%\pi^{\Lagr,-1}_{2d}(\cQ^{\sh})\cong \cB_{\cQ|_{Z^{\sh}}}\cong \cB_{\cQ^{\flat}}.
%\end{equation}
%This isomorphism restricts to an isomorphism $\th_{y'}: \cB^{\Herm}_{\cQ}(y')\isom \cB_{\cQ^{\flat}}(y)$ if $y'\in \Sig(Z^{\sh})$ maps to $y\in \Sig(Z)$. The map $\th_{y'}$ sends the chain $(\cQ_{\bu})$ to the non-strictly increasing chain of torsion sheaves $(\nu_{*}(\cQ_{\bu}|_{Z^{\sh}}))$.



%Let $d(Z)$ be the partition of $d$ given by $(d_{x})_{x\in Z}$. Fixing a total order on $Z$ makes the partition $d(Z)$ into a composition which we still denote by $d(Z)$, and gives a subgroup  $S_{d(Z)}\subset S_{d}$; taking preimage in $W_{d}$ gives a subgroup $W_{d(Z)}\cong \prod_{z\in Z}W_{d_{x}}\subset W_{d}$.
%
%
%\begin{cor}\label{c:Wind} There is an isomorphism of graded $W_{d}$-representations
%\begin{equation}
%\cohog{*}{\cB^{\SH}_{\cQ}}\cong \Ind^{W_{d}}_{W_{d(Z)}}\left(\bigotimes_{x\in Z}\cohog{*}{\cB^{\SH}_{\cQ_{\nu^{-1}(x)}}}\right).
%\end{equation}
%Here on the right side, each factor $W_{d_{x}}$ of $W_{d(Z)}$ acts on the tensor factor indexed by $x$ (for $x\in Z$).
%\end{cor}
%\begin{proof}
%The total order fixed on $Z$ gives a unique increasing function $y_{0}\in \Sigma(Z)$. Let $\Sigma_{0}\subset \Sigma(Z')$ be the preimage of $y_{0}$. For a subset $\Sigma'\subset \Sigma(Z')$, let $\cB^{\SH}_{\cQ}(\Sigma')$ be the disjoint union of $\cB^{\SH}_{\cQ}(y')$ for $y'\in \Sigma'$. Since $W_{d(Z)}$ stabilizes $\Sigma_{0}$, by Lemma \ref{l:permute comp}, $W_{d(Z)}$ stabilizes $\cohog{*}{\cB^{\SH}_{\cQ}(\Sigma_{0})}$. Moreover, $\cB^{\SH}_{\cQ}$ is the disjoint union of $\cB^{\SH}_{\cQ}(\Sigma_{0}\c w^{-1})$ for $w\in W_{d}/W_{d(Z)}$, and $\cohog{*}{\cB^{\SH}_{\cQ}(\Sigma_{0}\c w^{-1})}=w\cohog{*}{\cB^{\SH}_{\cQ}(\Sigma_{0})}$ by Lemma \ref{l:permute comp} again. We conclude that $\cohog{*}{\cB^{\SH}_{\cQ}}\cong \Ind^{W_{d}}_{W_{d(Z)}}\cohog{*}{\cB^{\SH}_{\cQ}(\Sigma_{0})}$. 
%
%We have a canonical isomorphism
%\begin{equation}
%\cB^{\SH}_{\cQ}(\Sigma_{0})\cong \prod_{x\in Z}\cB^{\SH}_{\cQ_{\nu^{-1}(x)}}.
%\end{equation}
%By K\"unneth formula we see that
%\begin{equation}\label{BS0}
%\cohog{*}{\cB^{\SH}_{\cQ}(\Sigma_{0})}\cong \ot_{x\in Z}\cohog{*}{\cB^{\SH}_{\cQ_{\nu^{-1}(x)}}}.
%\end{equation}
%To see that \eqref{BS0} is equivariant under $W_{d(Z)}$, the argument is the same as in the proof of Corollary \ref{c:Sind}, using a partial Springer resolution of $\SH_{d}(X'/X)$. \footnote{Y: More subtle than Coh. complete.}
%\end{proof}

%Choose a subset $Y\subset Z'$ such that $Y\coprod \s(Y)=Z'$. For each $z\in Z'$, let $\cB_{\cQ_{z}}$ be the type $A$ Springer fiber classifying full flags of coherent subsheaves $0\subset \cL_{1}\subset\cdots\subset \cL_{d_{z}}=\cQ_{z}$. Changing $z$ to $\s(z)$, we have a canonical isomorphism $\cB_{\cQ_{z}}\cong \cB_{\cQ_{\s(z)}}$ since $\cQ_{z}$ and $\cQ_{\s(z)}$ are dual to each other. 
%
%\begin{lemma}\label{l:SH fiber factor}
%For each $y'\in \Sigma(Z')$, there is a canonical isomorphism
%\begin{equation}
%\g_{y'}: \cB^{\SH}_{\cQ}(y')\cong \prod_{z\in Y}\cB_{\cQ_{z}}.
%\end{equation}
%\end{lemma}
%\begin{proof}
%For each $\cQ_{\bu}\in \cB^{\SH}_{\cQ}(y')$ and $z\in Y$, the distinct members of the sequence $\{\cQ_{i,z}\}$ give a full flag of coherent subsheaves of $\cQ_{z}$. This defines the projection $\cB^{\SH}_{\cQ}(y')\to \cB_{\cQ_{z}}$ for each $z\in Y$. Conversely, given a full flag $(\cL_{i}^{z})$ for each $\cQ_{z}$, we define $i(z)=\#\{1\le j\le i; y'(j)=z\}$  for each $z\in Z'$. For $z\in \s(Y)$ we define $\cL_{j}^{z}:=(\cL_{d_{z}-j}^{\s(z)})^{\bot}$ under the perfect pairing between $\cQ_{z}$ and $\cQ_{\s(z)}$. Then set $\cQ_{i}=(\bigoplus_{z\in Z'} \cL_{i(z)}^{z}$. This defines a full isotropic flag $(\cQ_{i})$ for $\cQ$ with $\supp(\cQ_{i}/\cQ_{i-1})=y'(i)$. It is easy to check that the two constructions are inverse to each other.
%\end{proof}



\subsection{Comparing Frobenius traces of $\HSpr_{d}$ and $\Spr_{d}$}\label{ss:trace comp Spr}
\source{higher-siegel-weil-unitary-nonsingular-terms}
In this subsection we will prove a relationship between Frobenius trace functions for $\HSpr_{d}$ and for $\Spr_{d}$. Since these sheaves live on different stacks, to make sense of the comparison we first need to identify the isomorphism classes of the $k$-points of these stacks.

For a groupoid $\cG$, let $|\cG|$ denote its set of isomorphism classes.
\begin{lemma}
\source{higher-siegel-weil-unitary-nonsingular-terms}
\notready\label{l:bij}
\source{higher-siegel-weil-unitary-nonsingular-terms}
There is a canonical bijection of sets
\begin{equation*}
|\Herm_{2d}(X'/X)(k)|\cong |\Coh_{d}(X)(k)|
%\cong |\SH_{2d}(X'/X)(k)|
\end{equation*}
respecting the support maps to $X_{d}(k)$. %\footnote{In the sequel we only use that there is a canonical map $|\Herm_{2d}(X'/X)(k)|\to |\Coh_{d}(X)(k)|$ respecting the support and Jordan types.}
\end{lemma}
\begin{proof} Let $\cP(d)$ be the set of partitions of $d\in\Z_{\ge0}$, and $\cP=\coprod_{d\ge0}\cP(d)$. Let $\cP_{|X|}$ be the set of functions $\l: |X|\to \cP$ such that $\l(v)$ is the zero partition for almost all $v\in |X|$. For $\l\in \cP_{|X|}$, let $|\l|=\sum_{v}|\l(v)|\deg(v)$. Let $\cP_{|X|}(d)$ be the subset of those $\l\in \cP_{|X|}$ with $|\l|=d$. Let $s_{d}: \cP_{|X|}(d)\to X_{d}(k)$ be the map sending  $\l$ to the divisor $\sum_{v}|\l(v)|v$. 

By taking the Jordan type of a torsion sheaf at each closed point, we get a canonical bijection $\L^{\Coh}_{d}: |\Coh_{d}(X)(k)|\isom\cP_{|X|}(d)$. The map $s_{d}^{\Coh}:|\Coh_d(X)(k)|\to X_{d}(k)$ corresponds to $s_{d}$ under this bijection.

We define a map $\L^{\Herm}_{d}: |\Herm_{2d}(X'/X)(k)|\to \cP_{|X|}(d)$ as follows. For $(\cQ,h)\in \Herm_{2d}(X'/X)(k)$ and $v\in |X|$,  let $\l(v)$ be the Jordan type of $\cQ_{v'}$ (the summand $\cQ_{v'}$ supported at $v'$) for any $v'\in |X'|$ above $v$. When $v$ is split in $X'$, the two choices of $v'$ give the same Jordan type. The support map $s_{2d}^{\Herm}: |\Herm_{2d}(X'/X)(k)|\to X_{d}(k)$ is the composition $s_{d}\c \L^{\Herm}_{d}$. 

We claim that $\L^{\Herm}_{d}$ is a bijection. Then $\L^{\Coh,-1}_{d}\c\L^{\Herm}_{d}: |\Herm_{2d}(X'/X)(k)|\isom |\Coh_{d}(X)(k)|$ is the desired bijection. 

To prove the claim, since any torsion coherent sheaf splits as a direct sum over the finitely many points in its support, it suffices to fix a closed point $v\in |X|$ and show that the set of isomorphism classes $|\Herm_{v}|$ of Hermitian torsion sheaves supported above $v$ maps bijectively (by the restriction of $\L^{\Herm}_{d}$) to $\cP$ which are supported at $v$. 

If $v$ splits into $v'$ and $v''$ in $|X'|$, then any $(\cQ,h)\in |\Herm_{v}|$ has the form $\cQ_{v'}\op \s^{*}\cQ_{v'}$ equipped with the canonical Hermitian structure. In this case we see that the isomorphism class of $(\cQ,h)$ is determined by the Jordan type of $\cQ_{v'}$, and conversely each Jordan type arises from some $(\cQ, h)$. 

If $v$ is inert with preimage  $v'\in |X'|$, let $(\cQ,h)\in |\Herm_{v}|$ be of length $d$ over $k_{v'}$. Then $V=\Gamma(X',\cQ)$ is a $d$-dimensional Hermitian $k_{v'}$-vector space with a self-adjoint nilpotent endomorphism $e$ given by the action of a uniformizer $\vp\in \cO_{v}$. Fix a $d$-dimensional Hermitian space $(V_{d},h)$ over $k_{v'}$ (unique up to isomorphism), then the isomorphism classes of $(\cQ,h)\in |\Herm_{v}|$ with length $d$ over $k_{v'}$ is in bijection with the adjoint orbits of the unitary group $\Ug(V_{d},h)(k_v)$ acting on the nilpotent cone $\cN(V_{d},h)(k_v)$ of self-adjoint nilpotent endomorphisms of $V_{d}$. Being a Galois twisted version of the usual nilpotent orbits under $\GL_{d}$, the orbits $\cN(V_{d},h)(k_v)/\Ug(V_{d},h)(k_v)$ are again classified by partitions of $d$ according the Jordan types of $e\in \cN(V_{d},h)(k_v)$ (here we use that the centralizer $C_{\GL_{d}}(e)$ is connected, and Lang's theorem implies $\cohog{1}{k_v, C_{\GL_{d}}(e)}=\{1\}$). Therefore the isomorphism class of $(\cQ,h)\in |\Herm_{v}|$ is determined by the Jordan type of $\cQ$, and conversely each Jordan type arises from a $(\cQ, h)$. This shows that $\L^{\Herm}_{d}$ is a bijection.  
% The argument for $\SH_{2d}(X'/X)$ is similar.
\end{proof}



\subsubsection{Further notations}\label{sss:choose points}
\source{higher-siegel-weil-unitary-nonsingular-terms}
Now let $(\cQ,h)\in \Herm_{2d}(k)$. We write $\cQ_{\ov k}$ for the base change of $\cQ$ over $X'_{\ov k}$, and adapt the notations $Z'\subset X'(\ov k),Z\subset X(\ov k),\Sig(Z'),\Sig(Z)$ from \S\ref{ss:Herm fiber}. Let $|Z'|$ and $|Z|$ be the set of closed points contained in $Z'$ and $Z$. We have a decomposition
\begin{equation*}
|Z|=|Z|_{s}\coprod |Z|_{i}
\end{equation*}
into split and inert places. For each closed point $v\in |Z|$ we choose a geometric point $x'_{v}\in Z'$ above $v$ and denote its image in $Z$ by $x_{v}$.

Let $\Frob:X'\to X'$ be the Frobenius morphism.  Let $Z^{\sh}$ be the following subset of $Z'$
\begin{equation*}
Z^{\sh}=\left\{F^{i}(x'_{v}): v\in |Z|, 0\le i<\deg(v)\right\}.
\end{equation*}
When $v$ splits into $v',v''$ in $|X'|$, with $x'_{v}|v'$, then $Z^{\sh}$ contains all geometric points above $v'$ and not any above $v''$. When $v$ is inert with preimage $v'\in|X'|$,  $Z^{\sh}$ contains half of the geometric points above $v'$ which form a chain under the Frobenius, starting with $x'_{v}$. Therefore $Z'=Z^{\sh}\coprod\s(Z^{\sh})$. For $x\in Z$ let $x^{\sh}\in Z^{\sh}$ be the unique element above $x$. This induces a section $\Sig(Z)\isom \Sig(Z^{\sh})\subset \Sig(Z')$ which we denote $y\mapsto y^{\sh}$.

Let $\cQ^{\flat}\in \Coh_{d}(X)(k)$ be the point corresponding to the isomorphism class of $(\cQ,h)$ under the bijection in  Lemma \ref{l:bij}. Then $\cQ^{\flat}_{\ov k}\cong\nu_{*}(\cQ|_{Z^{\sh}})$. Recall the isomorphism \eqref{BQy} for each $y\in \Sig(Z)$
\begin{equation*}
\b_{y}: \cB_{\cQ^{\flat}}(y)\isom \prod_{x\in Z}\cB_{\cQ^{\flat}_{x}}.
\end{equation*}
From this and the K\"unneth formula we get an identification
\begin{equation*}
\cohog{*}{\cB_{\cQ^{\flat}}(y)}=\bigotimes_{x\in Z} \cohog{*}{\cB_{\cQ^{\flat}_{x}}}.
\end{equation*}
By \cite[Corollary 2.3(1)]{HS77}, $\cohog{*}{\cB_{\cQ^{\flat}_{x}}}$ is concentrated in even degrees. Let $\cohog{*}{\cB_{\cQ^{\flat}}(y)}^{+}$ be the direct sum of all $\ot_{x\in Z} \cohog{2i_{x}}{\cB_{\cQ^{\flat}_{x}}}$ (for varying $(i_{x})_{x\in Z}$) such that 
\begin{equation}\label{sum ixv}
\source{higher-siegel-weil-unitary-nonsingular-terms}
\mbox{$\sum_{v\in|Z|_{i}}i_{x_{v}}$ is even.}
\end{equation}
Similarly, let $\cohog{*}{\cB_{\cQ^{\flat}}(y)}^{-}$ be the direct sum of all $\ot_{x\in Z} \cohog{2i_{x}}{\cB_{\cQ^{\flat}_{x}}}$ such that the  quantity in \eqref{sum ixv} is odd. We have
\begin{equation}\label{HBy pm}
\source{higher-siegel-weil-unitary-nonsingular-terms}
\cohog{*}{\cB_{\cQ^{\flat}}(y)}=\cohog{*}{\cB_{\cQ^{\flat}}(y)}^{+}\op\cohog{*}{\cB_{\cQ^{\flat}}(y)}^{-}.
\end{equation}
Taking direct sum over all $y\in \Sig(Z)$ we get a decomposition
\begin{equation}\label{HBQ pm}
\source{higher-siegel-weil-unitary-nonsingular-terms}
\cohog{*}{\cB_{\cQ^{\flat}}}=\cohog{*}{\cB_{\cQ^{\flat}}}^{+}\op\cohog{*}{\cB_{\cQ^{\flat}}}^{-}.
\end{equation}
Note that this decomposition depends on the choice of a geometric point $x_{v}$ over each inert $v$. By Corollary \ref{c:Sind}, the action of $S_{y}\subset S_{d}$ on $\cohog{*}{\cB_{\cQ^{\flat}}(y)}$ preserves the decomposition \eqref{HBy pm} since it is the same as the tensor product of the Springer actions on each factor $\cohog{*}{\cB_{\cQ^{\flat}_{x}}}$. Therefore the decomposition \eqref{HBQ pm} is stable under the $S_{d}$-action. 

Now $(\cQ|_{Z^{\sh}}\subset \cQ)$ gives a geometric point of $\Lagr^{\dm}_{2d}$, which is not defined over $k$ if $|Z|_{i}\ne\vn$.  Using this geometric point in $\Lagr^{\dm}_{2d}$, Proposition \ref{p:Spr HSpr} gives an isomorphism $\a^{\sh}:=\a_{(\cQ|_{Z^{\sh}}\subset\cQ)}$ on the level of stalks (see \eqref{astalk}):
\begin{equation*}
\a^{\sh}: \cohog{*}{\cB_{\cQ}^{\Herm}}^{(\Z/2\Z)^{d}}\cong \cohog{*}{\cB_{\cQ^{\flat}}}.
\end{equation*}
This isomorphism is $S_{d}$-equivariant. Both sides now carry geometric Frobenius actions which we denote by $\gFrob_{\cQ}$ and $\gFrob_{\cQ^{\flat}}$, which are not necessarily intertwined under $\a^{\sh}$ because the point $(\cQ|_{Z^{\sh}}\subset \cQ)$ is not necessarily defined over $k$. The next result gives the relation between the two Frobenius actions.


\begin{prop}
\source{higher-siegel-weil-unitary-nonsingular-terms}
\notready\label{p:trace relation}  Let $\th$ be the involution on $\cohog{*}{\cB_{\cQ^{\flat}}}$ which is $1$ on $\cohog{*}{\cB_{\cQ^{\flat}}}^{+}$ and $-1$ on $\cohog{*}{\cB_{\cQ^{\flat}}}^{-}$. Then under the isomorphism $\a^{\sh}$, $\gFrob_{\cQ}$ corresponds to $\gFrob_{\cQ^{\flat}}\c\th$.
\source{higher-siegel-weil-unitary-nonsingular-terms}
\end{prop}
\begin{proof} Recall from the proof of Proposition \ref{p:Spr HSpr} that $\a^{\sh}$ is the composition
\begin{equation*}
\cohog{*}{\cB_{\cQ}^{\Herm}}^{(\Z/2\Z)^{d}}\cong \bigoplus_{y\in \Sig(Z)}\cohog{*}{\cB^{\Herm}_{\cQ}(y^{\sh})}\xr{\op\g_{Z^{\sh}, y^{\sh}}}\bigoplus_{y\in \Sig(Z)}\cohog{*}{\cB_{\cQ^{\flat}}(y)}.
\end{equation*}
Here $\g_{Z^{\sh},y^{\sh}}$ is defined in \eqref{gZ}.


Then $\Fr^{*}$ on $\cohog{*}{\cB^{\Herm}_{\cQ}}$ maps $\cohog{*}{\cB^{\Herm}_{\cQ}(\Fr(y^{\sh}))}$ to $\cohog{*}{\cB^{\Herm}_{\cQ}(y^{\sh})}$. Note that $\Fr(y^{\sh})$ and $(\Fr y)^{\sh}$ are in general different: if $y(i)=\Fr^{-1}(x_{v})$ for some inert  $v$, then $\Fr(y^{\sh})(i)=\s(x'_{v})$ while $(\Fr y)^{\sh}(i)=x'_{v}$. In other words, the only difference between $\Fr (y^{\sh})$ and $(\Fr y)^{\sh}$ is the switch of all $x'_{v}$ and $\s(x'_{v})$ for all inert $v$. Therefore there is a unique element $\t_{y}\in(\Z/2\Z)^{d}$ such that $(\Fr y)^{\sh}=\Fr (y^{\sh})\c \t_{y}$.

Identifying $\cohog{*}{\cB_{\cQ}^{\Herm}}^{(\Z/2\Z)^{d}}$ with $\op_{y\in \Sig(Z)}\cohog{*}{\cB_{\cQ}^{\Herm}(y^{\sh})}$, the geometric Frobenius endomorphism $\gFrob_{\cQ}$ on $\cohog{*}{\cB_{\cQ}^{\Herm}}^{(\Z/2\Z)^{d}}$ is the direct sum of the following compositions
\begin{equation*}
\cohog{*}{\cB_{\cQ}^{\Herm}((\Fr y)^{\sh})}\xr{\t_{y}}\cohog{*}{\cB_{\cQ}^{\Herm}(\Fr (y^{\sh}))} \xr{\Fr^{*}}\cohog{*}{\cB_{\cQ}^{\Herm}(y^{\sh})}
\end{equation*}
where the first map is the Springer action of $\t_{y}\in W_{d}$ on $(\Spr^{\Herm}_{2d})_{\cQ}$.

On the other hand, let $w_{y}\in W_{d}$ be the {\em minimal length} element such that $(\Fr y)^{\sh}= \Fr(y^{\sh})\c w_{y}$. Write
\begin{equation*}
\t_{y}=w_{y}u_{y}, \mbox{ for a unique } u_{y}\in \Stab_{W_{d}}((\Fr y)^{\sh})=\Stab_{S_{d}}(\Fr y)\subset S_{d}.
\end{equation*}
Note that $\Stab_{S_{d}}(Fy)=\prod_{x\in Z}S_{I_{x}}$ where $I_{x}\subset \{1,2,\cdots, d\}$ is the preimage of $x$ under $y$. An easy calculation shows that $u_{y}=(u_{y,x})_{x\in Z}$ where  $u_{y,x}\in S_{I_{x}}$ is 
\begin{equation*}
u_{y,x}=\begin{cases} w_{I_{x}}, & \mbox{ if }x=x_{v}, v\in |Z|_{i},\\ 1, & \mbox{ otherwise.}\end{cases}
\end{equation*}
Here $w_{I_{x}}\in S_{I_{x}}$ is the involution that reverses the order of $I_{x}$. 


We use abbreviated notation
\begin{align*}
H(y') &:=\cohog{*}{\cB^{\Herm}_{\cQ}(y')}, \mbox{ for $y'\in\Sig(Z')$},\\
C(y) &:=\cohog{*}{\cB_{\cQ^{\flat}}(y)}, \mbox{ for $y\in\Sig(Z)$}.
\end{align*}
For each $y\in \Sig(Z)$, consider the following diagram
\begin{equation}\label{HC}
\source{higher-siegel-weil-unitary-nonsingular-terms}
\xymatrix{H((\Fr y)^{\sh})\ar[r]^{u_{y}}\ar[d]^{\g_{Z^{\sh}, (\Fr y)^{\sh}}} & H((\Fr y)^{\sh})\ar[rr]^-{w_{y}}\ar[d]^{\g_{Z^{\sh}, (\Fr y)^{\sh}}} & & H(\Fr (y^{\sh})) \ar[dll]^{\g_{Z^{\sh}, \Fr (y^{\sh})}}\ar[r]^{\Fr ^{*}}\ar[d]^{\g_{\Fr (Z^{\sh}), \Fr (y^{\sh})}} & H(y^{\sh})\ar[d]^{\g_{Z^{\sh}, y^{\sh}}}\\
C(Fy)\ar[r]^{u_{y}} & C(Fy)\ar[rr]^{\d^{*}} && C(Fy)\ar[r]^{\Fr^{*}} & C(y)}
\end{equation}
The left square is commutative by the $S_{d}$-equivariance of $\a^{\sh}$ proved in Proposition \ref{p:Spr HSpr} (here $u_{y}\in \Stab_{W_{d}}((\Fr y)^{\sh})\subset S_{d}$). The middle upper triangle is commutative by Lemma \ref{l:min w Herm}. The map $\d^{*}$ is defined to make the lower middle triangle commutative. The right square is clearly commutative.  The composite of the upper row is the restriction of $\Frob_{\cQ}$ to $H((\Fr y)^{\sh})$. Let us compute the composite of the lower row. 

The map $\d^{*}$ is the pullback along the automorphism $\d$ of $\cB_{\cQ^{\flat}}(\Fr y)$ that makes the following diagram commutative
\begin{equation*}
\xymatrix{& \cB^{\Herm}_{\cQ}(\Fr (y^{\sh})) \ar[d]^{\g_{\Fr(Z^{\sh}), \Fr(y^{\sh})}}\ar[dl]_{\g_{Z^{\sh}, \Fr(y^{\sh})}}\\
\cB_{\cQ^{\flat}}(\Fr y) & \cB_{\cQ^{\flat}}(\Fr y) \ar[l]_{\d}}
\end{equation*}
Under the isomorphism $\b_{\Fr y}:\cB_{\cQ^{\flat}}(\Fr y)\isom \prod_{x\in Z}\cB_{\cQ^{\flat}_{x}}$, $\d$ is the product of automorphisms $\d_{x}$ for each $\cB_{\cQ^{\flat}_{x}}$. If $x$ is not of the form $x=x_{v}$ for $v\in |Z|_{i}$, 
$\d_{x}$ is the identity. If $x=x_{v}$ for some $v\in |Z|_{i}$, then $\cQ^{\flat}_{x}=\cQ_{x'_{v}}$ and the Hermitian structure on $\cQ$ gives an isomorphism $\io: \cQ_{\s(x'_{v})}\cong\cQ_{x'_{v}}^{\vee}$. On the other hand, $\Fr^{\deg(v)}$ gives an isomorphism $\phi: \cQ_{x'_{v}}\cong \cQ_{\s(x'_{v})}$ since $\s(x'_{v})=\Fr^{\deg(v)}(x'_{v})$.  Combining $\io$ and $\phi$ we get a perfect symmetric pairing $(\cdot,\cdot)_{x'_{v}}$ on $\cQ_{x'_{v}}$ itself. Then $\d_{x}$ sends a full flag $\cR_{\bu}$ of $\cQ^{\flat}_{x_{v}}=\cQ_{x'_{v}}$ to $\cR^{\bot}_{\bu}$ under the pairing $(\cdot,\cdot)_{x'_{v}}$.

By the description of $u_{y}$ and $\d$ above, under the isomorphism $\b_{Fy}$, the composition $\d^{*}\c u_{y}$ takes the form
\begin{equation*}
\ot(\d_{x}^{*}\c u_{y,x}): \bigotimes_{x\in Z}\cohog{*}{\cB_{\cQ^{\flat}_{x}}}\to\bigotimes_{x\in Z}\cohog{*}{\cB_{\cQ^{\flat}_{x}}}.
\end{equation*}
The automorphisms $\d_{x}^{*}\c u_{y,x}$ are the identity maps except when $x=x_{v}$ for some $v\in |Z|_{i}$, in which case $u_{y,x}=w_{I_{x}}$. Let us compute $\d_{x}^{*}\c w_{I_{x}}$ on $\cohog{*}{\cB_{\cQ^{\flat}_{x}}}$ for $x=x_{v}$ and $v\in |Z|_{i}$.  For this we switch to the following notation. Let $V=\cQ_{x'_{v}}=\cQ^{\flat}_{x_{v}}$, a vector space of dimension $m$ over $\ov k$. We have argued that $V$ carries a symmetric self-duality $(\cdot,\cdot)$; the action of a uniformizer at $x_{v}$ gives a nilpotent element $e\in \End_{\ov k}(V)$, which is self-adjoint under $(\cdot,\cdot)$. Let $\cB$ be the flag variety of $\GL(V)$ and $\cB_{e}$ be the Springer fiber of $e$. Then $S_{m}$ acts on $\cohog{*}{\cB_{e}}$. Let $w_{0}$ be the longest element in $S_{m}$. Let $\d: \cB_{e}\isom \cB_{e}$ be the map sending a flag $V_{\bu}$ to $V_{\bu}^{\bot}$. We claim that $\d^{*}\c w_{0}$ acts on $\cohog{2i}{\cB_{e}}$ by $(-1)^{i}$. Indeed, by \cite[Corollary 2.3(2)]{HS77}, the restriction map $\cohog{*}{\cB}\to \cohog{*}{\cB_{e}}$ is surjective and is clearly equivariant under $\d^{*}\c w_{0}$, so it suffices to show that $\d^{*}\c w_{0}$ acts by $(-1)^{i}$ on $\cohog{2i}{\cB}$. Since $\d^{*}\c w_{0}$ preserves the cup product on $\cohog{*}{\cB}$, it suffices to show it acts by $-1$ on $\cohog{2}{\cB}$ (which generates all of $\cohog{2*}{\cB}$ under the cup product). For $1\le j\le m$, let $\xi_{j}$ be  the Chern class of the tautological line bundle on $\cB$ whose fiber at $V_{\bu}$ is  $V_j/V_{j-1}$. Then $\cohog{2}{\cB}$ is spanned by $\xi_j$ for $1\le j\le m$. Now we have $w_{0}(\xi_{j})=\xi_{m+1-j}$ since $\cohog{2}{\cB}$ is the reflection representation of $S_{m}$, and $\d^{*}\xi_{j}=-\xi_{m+1-j}$ by the definition of $\d$. Therefore $\d^{*}\c w_{0}(\xi_{j})=-\xi_{j}$ for all $1\le j\le m$, which proves that $\d^{*}\c w_{0}$ acts by $-1$ on $\cohog{2}{\cB}$, and hence acts by $(-1)^{i}$ on $\cohog{2i}{\cB}$ and on $\cohog{2i}{\cB_{e}}$. 

The above argument shows that $\d^{*}\c u_{y}$ acts by $1$ on $\cohog{*}{\cB_{\cQ^{\flat}}(\Fr y)}^{+}$ and acts by $-1$ on $\cohog{*}{\cB_{\cQ^{\flat}}(\Fr y)}^{-}$. Therefore the bottom row of \eqref{HC} is $\Frob_{\cQ^{\flat}}\c\th$. By the commutativity of \eqref{HC}, $\Frob_{\cQ^{\flat}}\c\th$ corresponds under $\a^{\sh}$ to the composite of the top row, which is $\Frob_{\cQ}$. This finishes the proof.
\end{proof}
